\documentclass[11pt]{article}
\usepackage[T1]{fontenc}
\usepackage{lmodern}
\usepackage[margin=1in]{geometry}
\usepackage{amsmath,amssymb,amsthm,mathtools}
\usepackage{microtype}
\usepackage[colorlinks=true,linkcolor=blue,citecolor=blue,urlcolor=blue]{hyperref}
\hypersetup{pdftitle={Smith normal form of the n-cycle class sum},pdfauthor={P. Graham}}
\newtheorem{theorem}{Theorem}[section]
\newtheorem{lemma}[theorem]{Lemma}
\newtheorem{corollary}[theorem]{Corollary}
\newtheorem{proposition}[theorem]{Proposition}
\theoremstyle{remark}
\newtheorem{remark}[theorem]{Remark}
\newcommand{\Z}{\mathbb{Z}}
\newcommand{\Q}{\mathbb{Q}}
\newcommand{\im}{\operatorname{im}}
\newcommand{\diag}{\operatorname{diag}}
\newcommand{\SNF}{\operatorname{SNF}}
\title{Smith normal form of the n-cycle class sum}
\author{P. Graham\\ \small\url{https://orcid.org/0009-0007-8493-1505}}
\date{October 22, 2026}
\begin{document}
\maketitle
\begin{abstract}
Let $M_n$ be the integer matrix of multiplication by the sum of all
$n$-cycles on the center of $\Z S_n$, in the conjugacy-class-sum basis.
We determine its Smith normal form for every positive integer $n$.
It has $p(n)-n$ zero entries and nonzero invariant factors $d_0=1$ and
$d_r=a_r r!/\gcd(n,r)$ for $1\leq r<n$, where $a_r=2$ when
$r\equiv n\pmod 2$ and $a_r=1$ otherwise, and $p(n)$ is the partition number.
This confirms Stanley's two conjectures for odd primes and twice odd primes,
posed in 2019, when he also suggested that further data might lead to a
conjecture for all $n$.
Starting from the classical hook-character expression for the structure
constants, we determine two integral coefficient lattices and reduce the
induced pairing to an anti-triangular matrix with rowwise divisibility.
The image-lattice theorem and the divisor chain are formalized in Lean
from group-ring definitions, including a proof of the character expression.
The equivalent diagonal form and the prime cases follow by the written
basis-change and specialization arguments.
\end{abstract}

\noindent Checked in Lean from the definitions; not yet reviewed by experts.

\smallskip
\noindent The proofs, the Lean formalization and the text were produced by AI models
(from OpenAI and Anthropic), working under the author's direction.

\section{Introduction}
For a partition $\mu\vdash n$, write $K_\mu$ for the sum of the
permutations of cycle type $\mu$ in $\Z S_n$. The elements $K_\mu$ form
a $\Z$-basis of the center $Z(\Z S_n)$, identifying it with
$\Z^{p(n)}$, where $p(n)$ is the number of partitions of $n$.
We consider the matrix $M_n$ defined by
\begin{equation}\label{eq:matrix}
 K_{(n)}K_\nu=\sum_{\mu\vdash n}M_n[\mu,\nu]K_\mu.
\end{equation}
Thus the rows record output cycle types and the columns record input cycle
types. The Smith normal form of an integer matrix $A$ is a diagonal matrix
obtained by invertible integral row and column operations, with positive
nonzero entries $d_0\mid d_1\mid\cdots$ followed by zeros. Equivalently,
it describes the embedding of the column lattice of $A$ in its ambient
free abelian group up to integral changes of basis.

Stanley asked for the Smith normal forms of class-sum multiplication
matrices in 2019~\cite{StanleyQuestion}, inspired by Lindzey's question
about the Cayley graph generated by $n$-cycles~\cite{Lindzey}.
Lindzey considers the full $n!$-dimensional group algebra; here the
operator is restricted to its center. For the $n$-cycle class Stanley gave
data, observed that the rank is $n$, and conjectured explicit formulas
when $n$ is an odd prime and when it is twice an odd prime. He suggested
that more data might lead to a conjecture for all $n$. His broader
question concerns arbitrary conjugacy classes and finite groups.
The following theorem determines the $n$-cycle case for every $n$.

\begin{theorem}\label{thm:main}
Let $n\geq1$, and let $p(n)$ denote the number of partitions of $n$.
Define
\begin{equation}\label{eq:d}
 d_0=1,\qquad
 d_r=\frac{a_r r!}{\gcd(n,r)}\quad(1\leq r<n),\qquad
 a_r=\begin{cases}2&r\equiv n\pmod2,\\1&r\not\equiv n\pmod2.\end{cases}
\end{equation}
These are positive integers satisfying $d_r\mid d_{r+1}$ for
$0\leq r<n-1$. There is a basis
\[
 b_0,\ldots,b_{n-1},b'_1,\ldots,b'_{p(n)-n}
\]
of $\Z^{p(n)}$ such that
\begin{equation}\label{eq:image}
 M_n\Z^{p(n)}=\bigoplus_{r=0}^{n-1}d_r\Z b_r.
\end{equation}
Consequently,
$\SNF(M_n)=\diag(d_0,\ldots,d_{n-1},0,\ldots,0)$,
with $p(n)-n$ zeros.
\end{theorem}

For an odd prime $n=p$, the factors are $r!$ at even indices and $2r!$
at odd indices; for $n=2p$, they are the factorials with $p!$ replaced
by $(p-1)!$. Thus both of Stanley's conjectures follow
(Corollary~\ref{cor:primes}). Section~\ref{sec:examples} compares the
formula with his examples, including a discrepancy at index six in
the displayed $n=12$ list.

The proof starts with the hook-character generating polynomials in $q$,
substitutes $q=-x$, and passes to coefficients in the shifted variable
$y=x-1$.
We determine the integral lattice spanned by these coefficient vectors
and the lattice obtained by dividing each vector by the corresponding
centralizer order. In these coordinates the classical character pairing
is anti-triangular. Bases for the two lattices expose its anti-diagonal
entries, and row divisibility then identifies the Smith factors.

The character-theoretic study of factorizations into $n$-cycles includes
Stanley's work~\cite{StanleyFactorizations}. Goupil and
Schaeffer~\cite{GoupilSchaeffer} give an explicit formula for factoring
an $n$-cycle into two permutations with prescribed cycle types, and
Biane~\cite{Biane} gives an algebraic proof of that formula.
These connection coefficients are related to the entries of $M_n$ by
the normalization stated in Section~\ref{sec:characters}.
The present proof determines the integral image lattice of this matrix.

Stanley's survey~\cite{StanleySurvey} gives the broader context of
Smith forms in combinatorics. Cai and Stanley~\cite{CaiStanley}
determine the Smith form of the up--down operator on Young's lattice,
equivalently a differential operator on symmetric functions. That
operator provides related context for~\eqref{eq:matrix}. The representation-theoretic
ingredients below are classical; see, for example, Sagan~\cite{Sagan}.
Grinberg's recent exposition~\cite{Grinberg} provides a detailed treatment
of class sums and symmetric-group representations, and was a source
for this background.

The Lean development proves~\eqref{eq:image} and the divisor chain from
the definitions of permutations and class sums, including a proof of
the character formula. The passage from
\eqref{eq:image} to diagonal Smith form and the prime specializations
are written deductions given below. Section~\ref{sec:formal} states the
formal theorem precisely and explains its connection with the written
proof.

\section{The character formula and a change of basis}\label{sec:characters}
Fix $n\geq1$ and put $N=n-1$. For $\mu\vdash n$, let $\ell(\mu)$ be
its number of parts and $m_j(\mu)$ the multiplicity of $j$. Set
\[
 z_\mu=\prod_{j\geq1}j^{m_j(\mu)}m_j(\mu)!,\qquad
 \epsilon_\mu=(-1)^{\ell(\mu)-1}.
\]
The number $z_\mu$ is the order of the centralizer of a permutation of
type $\mu$. Write $\chi_k=\chi^{(n-k,1^k)}$ for the hook character,
where $0\leq k\leq N$. The hook characters are the characters of the
exterior powers of the standard representation, so
\begin{equation}\label{eq:hook}
 \sum_{k=0}^N\chi_k(\mu)q^k
 =\frac{\prod_{j=1}^{\ell(\mu)}(1-(-q)^{\mu_j})}{1+q}.
\end{equation}
Indeed, the characteristic polynomial on the permutation representation
is the product of the characteristic polynomials of its cycles. Removing
its one-dimensional trivial summand divides the exterior-power generating
function by $1+q$.

\begin{proposition}\label{prop:bridge}
For $\mu,\nu\vdash n$,
\begin{equation}\label{eq:bridge}
 M_n[\mu,\nu]=\frac1{z_\nu}\sum_{k=0}^N
       (-1)^k k!(N-k)!\chi_k(\mu)\chi_k(\nu).
\end{equation}
\end{proposition}
\begin{proof}
Frobenius's class-multiplication formula gives
\[
 M_n[\mu,\nu]=\frac{|K_{(n)}|\,|K_\nu|}{n!}
 \sum_{\lambda\vdash n}
 \frac{\chi^\lambda((n))\chi^\lambda(\nu)\chi^\lambda(\mu)}{\chi^\lambda(1)}.
\]
All symmetric-group characters here are real, and every conjugacy class
is closed under inversion. By the Murnaghan--Nakayama rule, only hooks
contribute: $\chi_k((n))=(-1)^k$. Their dimensions are
$\binom Nk$, while $|K_{(n)}|=N!$ and $|K_\nu|=n!/z_\nu$.
Substitution gives~\eqref{eq:bridge}.
\end{proof}

To relate the normalization to the factorization literature, let
$c^{(n)}_{\mu\nu}$ count the factorizations of a fixed $n$-cycle as
a product of permutations of types $\mu$ and $\nu$. Equation~(2) of
Biane~\cite{Biane}, identified there with equation~(4) of
Goupil and Schaeffer~\cite{GoupilSchaeffer}, gives
\[
 c^{(n)}_{\mu\nu}=\frac{n}{z_\mu z_\nu}
 \sum_{k=0}^N(-1)^k k!(N-k)!\chi_k(\mu)\chi_k(\nu).
\]
Thus $M_n[\mu,\nu]=(z_\mu/n)c^{(n)}_{\mu\nu}$.
The multiplier depends on the row and can be nonintegral; this
normalization identity is not an equivalence under unimodular operations.
The character expression supplies the entries, while the lattice
arguments below determine their integral invariant factors.

Introduce polynomials
\begin{align*}
 H_\mu(y)&=\prod_{j=1}^{\ell(\mu)}((1+y)^{\mu_j}-1),&
 \eta_\mu(i)&=[y^{i+1}]H_\mu(y),\quad 0\leq i\leq N,\\
 \Phi_\mu(y)&=\frac{H_\mu(y)}{z_\mu},&
 e_s&=s!\gcd(n,s+1),\quad 0\leq s\leq N.
\end{align*}
The order of $H_\mu$ is $\ell(\mu)$, and its lowest coefficient is
$\prod_j\mu_j$. For $0\leq i,s\leq N$, define
\begin{equation}\label{eq:Q}
 Q(i,s)=\begin{cases}
 (-1)^{i+s}\dfrac{(N-i)!(N-s)!}{(N-i-s)!},&i+s\leq N,\\
 0,&i+s>N.
 \end{cases}
\end{equation}

\begin{lemma}\label{lem:pairing}
For $0\leq i,s\leq N$,
\[
 \sum_{k=0}^N(-1)^k k!(N-k)!
       [x^k](x-1)^i\,[x^k](x-1)^s=Q(i,s).
\]
Consequently,
\begin{equation}\label{eq:changed}
 M_n[\mu,\nu]=\frac{\epsilon_\mu\epsilon_\nu}{z_\nu}
        \sum_{i,s=0}^N\eta_\mu(i)Q(i,s)\eta_\nu(s).
\end{equation}
\end{lemma}
\begin{proof}
The coefficient $[x^k](x-1)^i$ is $(-1)^{i-k}\binom ik$.
Using
$k!\binom ik(N-k)!=i!(N-i)!\binom{N-k}{N-i}$,
the sum in the lemma becomes
\[
 (-1)^{i+s}i!(N-i)!
 \sum_{k=0}^s(-1)^k\binom sk\binom{N-k}{N-i}.
\]
Pascal's identity, applied successively $s$ times, gives
\[
 \sum_{k=0}^s(-1)^k\binom sk\binom{N-k}{m}
       =\binom{N-s}{m-s},
\]
where an out-of-range lower index means zero. Take $m=N-i$ and
simplify factorials to obtain~\eqref{eq:Q}.
Now set $g_\mu(x)=\prod_j(1-x^{\mu_j})/(1-x)$.
Equation~\eqref{eq:hook} implies $\chi_k(\mu)=(-1)^k[x^k]g_\mu(x)$,
and substitution gives
\[
 g_\mu(1+y)=\epsilon_\mu H_\mu(y)/y
            =\epsilon_\mu\sum_i\eta_\mu(i)y^i.
\]
Apply the first assertion to~\eqref{eq:bridge}.
\end{proof}

\section{The normalized coefficient lattice}
Define the rational vector $\alpha(\nu)$ by
\[
 \alpha(\nu)_s=\epsilon_\nu e_s[y^{s+1}]\Phi_\nu
       =\frac{\epsilon_\nu e_s\eta_\nu(s)}{z_\nu}.
\]

\begin{lemma}\label{lem:A}
The vectors $\alpha(\nu)$, for $\nu\vdash n$, belong to $\Z^n$ and
span $\Z^n$ over $\Z$.
\end{lemma}
\begin{proof}
We first prove integrality. Allow $\mu$ to be any multiset of positive
integers, including the empty multiset, and put
\[
 \phi_j(y)=\frac{(1+y)^j-1}{j},\qquad
 \Phi_\mu=\prod_{j\geq1}\frac{\phi_j^{m_j(\mu)}}{m_j(\mu)!},
 \qquad I_t(\mu)=t![y^t]\Phi_\mu.
\]
For the empty multiset, $\Phi_\varnothing=1$. Removing one part $j$
is denoted by $\mu-j$. Since $\phi_j'=(1+y)^{j-1}$,
\begin{equation}\label{eq:derivative}
 \Phi_\mu'=\sum_{j:m_j(\mu)>0}(1+y)^{j-1}\Phi_{\mu-j}.
\end{equation}
Induction on the sum of the parts proves $I_t(\mu)\in\Z$ for every
$t\geq0$. The constant coefficient is immediate; for $t\geq1$,
\eqref{eq:derivative} expresses $I_t(\mu)$ as a sum of
\[
 \binom{j-1}{i}\frac{(t-1)!}{(t-1-i)!}I_{t-1-i}(\mu-j),
 \qquad 0\leq i\leq\min(j-1,t-1),
\]
each an integer by induction.

When the sum of the parts is $n$, the identity
$j m_j(\mu)\Phi_\mu=((1+y)^j-1)\Phi_{\mu-j}$ yields
\begin{equation}\label{eq:euler}
 n\Phi_\mu=\sum_{j:m_j(\mu)>0}((1+y)^j-1)\Phi_{\mu-j}.
\end{equation}
For $t\geq1$, multiply the coefficient of $y^t$ by $(t-1)!$.
The right side is a sum of
\[
 \binom ji\frac{(t-1)!}{(t-i)!}I_{t-i}(\mu-j),
 \qquad 1\leq i\leq\min(j,t),
\]
so $nI_t(\mu)/t$ is an integer. B\'ezout's identity for $n$ and $t$
then shows that $\gcd(n,t)I_t(\mu)/t$ is an integer. Setting $t=s+1$
proves the required integrality.

For spanning, fix a length $\ell=s+1$. The vectors from partitions
of this length vanish in coordinates below $s$, and their $s$th
coordinates, up to sign, are
\begin{equation}\label{eq:leading-alpha}
 \frac{(\ell-1)!\gcd(n,\ell)}{\prod_j m_j(\nu)!}.
\end{equation}
We show that the gcd of these integers is one. If $\ell=n$, the partition
$1^n$ gives one. Suppose $\ell<n$, fix a prime $p$, and write
$v=v_p(\gcd(n,\ell))$. If $v=0$, the partition
$(n-\ell+1,1^{\ell-1})$ gives $\gcd(n,\ell)$, prime to $p$.
If $v>0$, put $q=p^v$ and $A=(n-\ell)/q+1\geq2$. The partition
$(A^q,1^{\ell-q})$ gives
\[
 \binom{\ell-1}{q-1}\frac{\gcd(n,\ell)}q.
\]
The second factor is prime to $p$. In
$\binom{\ell-1}{q-1}=\prod_{i=1}^{q-1}(\ell-i)/i$,
the valuations of numerator and denominator agree term by term, because
$v_p(i)<v\leq v_p(\ell)$. Thus the first factor is also prime to $p$.
No prime divides every number in~\eqref{eq:leading-alpha}.
B\'ezout therefore gives a combination $f_s$ of the length-$s+1$
vectors with coordinate one at $s$ and zero below it. The vectors
$f_0,\ldots,f_N$ form a unitriangular integral matrix, hence a basis
of $\Z^n$.
\end{proof}

\section{The integral coefficient lattice}
Let $U$ be the $\Z$-span of the vectors $\eta_\mu$ in $\Z^n$.
Extend the notation~\eqref{eq:d} by setting $a_0=n$, so that
$d_r=a_rr!/\gcd(n,r)$ also holds at $r=0$, using $\gcd(n,0)=n$.

\begin{lemma}\label{lem:B}
The lattice $U$ has a basis $u_0,\ldots,u_N$ such that $u_r(i)=0$
for $i<r$ and $u_r(r)=a_r$. Each $u_r$ is an integral combination of
the $\eta_\mu$ with $\ell(\mu)=r+1$.
\end{lemma}
\begin{proof}
Take $u_0=\eta_{(n)}$. When $N>0$, take $u_N=\eta_{1^n}$.
For $1\leq r<N$, the two partitions
\[
 (n-r,1^r),\qquad(n-r-1,2,1^{r-1})
\]
have length $r+1$ and leading coefficients $n-r$ and $2(n-r-1)$.
Their gcd is $\gcd(n-r,2)=a_r$, so B\'ezout gives the desired $u_r$.
This constructs independent vectors with the asserted leading terms.

We prove that they span. Identify $v\in\Z^n$ with
$\sum_{i=0}^Nv(i)y^i$. Every $v\in U$ satisfies $n\mid v(0)$.
It also satisfies
\begin{equation}\label{eq:parity}
 T_N(v)\equiv v\pmod2,\qquad
 T_m(f)=(1+y)^m f\bigl(y/(1+y)\bigr).
\end{equation}
Indeed, $T_{a+b}(fg)=T_a(f)T_b(g)$ and
\[
 T_m((1+y)^m-1)=(1+2y)^m-(1+y)^m
                  \equiv(1+y)^m-1\pmod2.
\]
Consequently $T_n(H_\mu)\equiv H_\mu\pmod2$.
Since $T_n(H_\mu)=yT_N(H_\mu/y)$, cancellation of $y$ proves
\eqref{eq:parity} on the generators and hence on $U$.

Suppose $v(i)=0$ for $i<r$. At $r=0$ we have $a_0\mid v(0)$.
For $1\leq r<N$, comparing coefficients of $y^{r+1}$ in
\eqref{eq:parity} gives $(N-r)v(r)\equiv0\pmod2$.
When $r\equiv n\pmod2$, the factor $N-r$ is odd, so $a_r=2$
divides $v(r)$. In all other cases, including $r=N$, $a_r=1$.
Subtracting $(v(r)/a_r)u_r$, successively for $r=0,\ldots,N$,
therefore expresses every element of $U$ in the constructed basis.
\end{proof}

\section{Divisibility of the pairing}
For a partition $\mu$ and $0\leq s\leq N$, put
\[
 \Gamma(\mu,s)=\frac1{e_s}\sum_{i=0}^N\eta_\mu(i)Q(i,s).
\]

\begin{lemma}\label{lem:C}
If $\ell(\mu)=r+1$, then $\Gamma(\mu,s)\in d_r\Z$ for every $s$.
\end{lemma}
\begin{proof}
Rewrite~\eqref{eq:Q} as
\[
 Q(i,s)=(-1)^{i+s}s!(N-s)!\binom{N-i}s.
\]
Thus, for the integer
$X_s=\sum_i(-1)^i\binom{N-i}s\eta_\mu(i)$,
\begin{equation}\label{eq:gamma}
 \Gamma(\mu,s)=(-1)^s\frac{(N-s)!}{\gcd(n,s+1)}X_s.
\end{equation}
Terms vanish unless $r\leq i\leq N-s$. If $s>N-r$ the assertion
follows at once. If $s=N-r$, then $X_s=(-1)^r\prod_j\mu_j$ and
\[
 \Gamma(\mu,N-r)=\pm\frac{r!}{\gcd(n,r)}\prod_j\mu_j.
\]
The product is divisible by $a_r$: for $r=0$ it equals $n$;
for $a_r=2$ and $r\geq1$, $r+1$ odd parts would have sum of
parity $r+1$, contrary to $n\equiv r\pmod2$.

It remains to treat $s<N-r$. Put $t=N-s>r$ and $b=\gcd(n,s+1)$.
As $b\mid n-(s+1)=t$, we have
\[
 \frac{t!}{b}=r!R,\qquad
 R=(r+1)(r+2)\cdots(t-1)\frac tb\in\Z.
\]
It suffices to show $a_r\mid\gcd(n,r)R$. This is immediate if
$r=0$, or $a_r=1$, or $n,r$ are even. The remaining case has
$a_r=2$ and $n,r$ odd. Then $b$ is odd. If $t>r+1$, $R$ contains
the even factor $r+1$; if $t=r+1$, its factor $t/b$ is even.
This proves the required divisibility, including integrality.
\end{proof}

\section{Assembly and the Smith form}
For $0\leq r,s\leq N$, define
\[
 G(r,s)=\frac1{e_s}\sum_i u_r(i)Q(i,s).
\]
Lemma~\ref{lem:C} and the length restriction in Lemma~\ref{lem:B}
show that $d_r$ divides every entry of row $r$. Moreover $G(r,s)=0$
if $r+s>N$, and
\[
 G(r,N-r)=\frac{a_rQ(r,N-r)}{e_{N-r}}
         =(-1)^N\frac{a_rr!}{\gcd(n,n-r)}=(-1)^Nd_r.
\]
Thus
\begin{equation}\label{eq:factorG}
 G=\diag(d_0,\ldots,d_N)V,
\end{equation}
where $V$ is an integral anti-triangular matrix with entries $\pm1$
on the anti-diagonal. Its determinant is $\pm1$, so $V$ is unimodular.

Write $\eta_\mu=\sum_r X(\mu,r)u_r$ using Lemma~\ref{lem:B}, and
write $u_r=\sum_\mu W(r,\mu)\eta_\mu$ using the construction there.
The matrices $X$ and $W$ are integral, and independence of the $u_r$
gives $WX=I_n$. Let $E=\diag(\epsilon_\mu)$ and let $\alpha$ be the
matrix whose columns are $\alpha(\nu)$. Equation~\eqref{eq:changed}
and the definitions give
\[
 M_n=E\Gamma\alpha=EXG\alpha.
\]
By Lemma~\ref{lem:A}, $\alpha:\Z^{p(n)}\to\Z^n$ is surjective.
Since $V$ is unimodular,~\eqref{eq:factorG} now yields
\[
 M_n\Z^{p(n)}=EX\diag(d_0,\ldots,d_N)\Z^n
             =\bigoplus_{r=0}^N d_r\Z b_r,
 \qquad b_r=EX\mathbf e_r.
\]
The map $WE$ is a left inverse of $EX$, because $E^2=I$.
Therefore
\[
 \Z^{p(n)}=EX\Z^n\oplus\ker(WE).
\]
The kernel is a free abelian group, being a subgroup of a finitely
generated free abelian group. Choosing its basis extends the $b_r$
to a basis of $\Z^{p(n)}$, proving~\eqref{eq:image}.

For completeness, all $d_r$ are positive integers: $d_0=1$, and for
$r\geq1$ the divisor $\gcd(n,r)$ divides $r!$. We also have
$d_0\mid d_1$ whenever the latter occurs. For $1\leq r<N$,
\[
 \frac{d_{r+1}}{d_r}
 =\frac{a_{r+1}}{a_r}\frac{(r+1)\gcd(n,r)}{\gcd(n,r+1)}.
\]
Since $\gcd(n,r+1)\mid r+1$, only $a_r=2$ could cause a
denominator. In that case $a_{r+1}=1$. If $n,r$ are even,
$\gcd(n,r)$ is even and $\gcd(n,r+1)$ is odd. If $n,r$ are odd,
$r+1$ is even and $\gcd(n,r+1)$ is odd. The ratio is integral
in both cases, establishing the divisor chain.

Finally, choose an integral change of basis $P$ carrying $b_r$ to
the first $n$ standard basis vectors. Then $PM_n$ maps
$\Z^{p(n)}$ onto $\bigoplus_r d_r\Z\mathbf e_r$. Lift the basis
vectors $d_r\mathbf e_r$ of this image to the domain. These lifts,
together with a basis of the kernel, form a basis of the domain:
the map onto its free image splits. The corresponding unimodular
column change $Q$ satisfies
\[
 PM_nQ=\diag(d_0,\ldots,d_{n-1},0,\ldots,0).
\]
The image has rank $n$, so precisely $p(n)-n$ entries vanish.
This proves Theorem~\ref{thm:main}.

\section{Stanley's conjectures and examples}\label{sec:examples}
\begin{corollary}\label{cor:primes}
If $n=p$ is an odd prime, the nonzero invariant factor at index $r$
is $r!$ for even $r$ and $2r!$ for odd $r$, for $0\leq r<p$.
If $n=2p$ for an odd prime $p$, the list is $0!,1!,\ldots,(2p-1)!$
with $p!$ replaced by $(p-1)!$.
\end{corollary}
\begin{proof}
For $1\leq r<p$, $\gcd(p,r)=1$, and $r\equiv p\pmod2$
means $r$ is odd. This proves the first assertion, including $r=0$
separately. When $n=2p$, even $r$ with $0<r<2p$ have gcd two
and $a_r=2$. Odd $r\ne p$ have gcd one and $a_r=1$. At $r=p$
the factor is $p!/p=(p-1)!$. Again $d_0=1$.
\end{proof}

The theorem also gives
\[
 Z(\Z S_n)/K_{(n)}Z(\Z S_n)
 \cong\Z^{p(n)-n}\oplus\bigoplus_{r=0}^{n-1}\Z/d_r\Z.
\]
These corollaries are immediate written consequences of the lattice
statement; separate Lean declarations for them are not needed for the
formal theorem in Section~\ref{sec:formal}.

For $n=1$, the matrix is $[1]$. The nonzero factors for $n=2,3,4,5$
are respectively
\[
 (1,1),\quad(1,2,2),\quad(1,1,2,6),\quad(1,2,2,12,24).
\]
The formula at $n=9$ gives
\[
 0!,\ 2\cdot1!,\ 2!,\ \tfrac23\,3!,\ 4!,\ 2\cdot5!,\
 \tfrac13\,6!,\ 2\cdot7!,\ 8!,
\]
agreeing term by term with Stanley's example~\cite{StanleyQuestion}.
At $n=12$, it gives
\[
 0!,\ 1!,\ 2!,\ \tfrac13\,3!,\ \tfrac12\,4!,\ 5!,\
 \tfrac13\,6!,\ 7!,\ \tfrac12\,8!,\ \tfrac13\,9!,\ 10!,\ 11!.
\]
The entry at index six is $6!/3$. The displayed $n=12$ example in
\cite{StanleyQuestion} instead has $2\cdot6!$ at this position;
that entry does not divide the following $7!$ and appears to be a
typographical error. The other entries agree. These examples illustrate
the formula and are not inputs to its proof.

\section{Formalization and reproducibility}\label{sec:formal}
The accompanying Lean project defines $S_n$ as the permutations of
$\{0,\ldots,n-1\}$ and works in its integral monoid algebra.
The class sum includes each permutation of the prescribed cycle type
with coefficient one. The matrix entry is the coefficient of a chosen
permutation of type $\mu$ in $K_{(n)}K_\nu$; the development proves
that the choice has the required cycle type and that these entries
give the expansion~\eqref{eq:matrix}.

The final theorem is:
\begin{verbatim}
theorem stanley_class_sum_snf (n : Nat) (hn : 0 < n) :
    HasStanleySNF n (M n) \u2227
      (\u2200 r, r + 1 < n \u2192 stanleyD n r \u2223 stanleyD n (r + 1))
\end{verbatim}
Here \texttt{stanleyD} is~\eqref{eq:d}, with its division proved exact.
The predicate \texttt{HasStanleySNF} asserts the existence of a finite
index set $\kappa$ and an integral basis indexed by
$\{0,\ldots,n-1\}\sqcup\kappa$, for which the image is the span of
$d_r b_r$. This is precisely~\eqref{eq:image}. Positivity and
independence turn that span into the direct sum stated in the theorem.
The written basis-change argument in Section~6 identifies it with the
usual diagonal Smith form. In particular the formal statement does not
replace the group-ring matrix by a matrix defined through a character
formula.

For the character formula, the formal proof uses an explicit model
of the exterior algebra of $\Q^n$, with basis indexed by subsets.
Permutations act by signed permutations of these basis vectors.
Exterior multiplication by $e=e_1+\cdots+e_n$ and contraction by
the all-ones linear functional give operators $E$ and $I$ satisfying
$IE+EI=n$. Thus $P=IE/n$ is an equivariant projection onto
$\ker I$. Taking traces of the degree-$k$ action on this image gives
the coefficients in~\eqref{eq:hook}.

The degree-$k$ commutant is determined by the signed action on pairs
of $k$-subsets: a transposition in their difference kills an entry when
their intersection has size at most $k-2$, and the remaining orbits
leave a two-dimensional span of the identity and the projection.
Restriction to the image of $P$ makes every commuting operator scalar.
Taking traces of conjugacy averages consequently gives
\[
 \sum_{g\in S_n}\chi_k(agb g^{-1})
       =\frac{n!}{\binom Nk}\chi_k(a)\chi_k(b).
\]
Evaluation of~\eqref{eq:hook} at $q=-1$ gives the indicator identity
$\sum_k(-1)^k\chi_k(\mu)=n\,\mathbf1_{\mu=(n)}$.
This identity, the preceding conjugacy average, and orbit--stabilizer
prove Proposition~\ref{prop:bridge} directly in Lean. Consequently,
the formal proof does not take the classical representation-theoretic
theorems used in its shorter written derivation as extra hypotheses.

The remaining development follows the polynomial pairing, the two
lattice lemmas, row divisibility, and assembly above. It uses Lean's
kernel and the mathlib library; no certificate generator, numerical
approximation, or external computational oracle is required. From the
bundle's \texttt{lean/} directory, the standard commands are
\begin{verbatim}
lake exe cache get
lake build
\end{verbatim}
The build includes checks of the final statement and its logical
dependencies. The accompanying Lean guide documents these checks and
the correspondence of definitions. The paper can be rebuilt from its
source with \texttt{pdflatex}.

\section*{Acknowledgments}
The question, the two prime-case conjectures, and the motivating
examples are due to Richard P. Stanley. The proof uses classical
character theory, as presented in the works of Sagan and Grinberg,
and the mathematical infrastructure provided by Lean and mathlib.
Cai and Stanley's work is credited as related
background, not as a prior proof for the operator considered here.

\begin{thebibliography}{9}
\bibitem{Biane}
P. Biane, \emph{Nombre de factorisations d'un grand cycle},
S\'eminaire Lotharingien de Combinatoire \textbf{51} (2004), Article B51a.
\url{https://www.mat.univie.ac.at/~slc/wpapers/s51biane.pdf}.

\bibitem{CaiStanley}
T. W. Cai and R. P. Stanley,
\emph{The Smith Normal Form of a Matrix Associated with Young's Lattice},
2015. \href{https://arxiv.org/abs/1502.00922}{arXiv:1502.00922}.

\bibitem{GoupilSchaeffer}
A. Goupil and G. Schaeffer,
\emph{Factoring $n$-cycles and counting maps of given genus},
European Journal of Combinatorics \textbf{19} (1998), no.~7, 819--834.
\href{https://doi.org/10.1006/eujc.1998.0215}{doi:10.1006/eujc.1998.0215}.

\bibitem{Grinberg}
D. Grinberg, \emph{An introduction to the symmetric group algebra}, 2025.
\href{https://arxiv.org/abs/2507.20706}{arXiv:2507.20706}.

\bibitem{Lindzey}
N. Lindzey, \emph{Smith Normal Form of a Cayley Graph of the Symmetric Group},
MathOverflow, June 27, 2019.
\url{https://mathoverflow.net/questions/334899}.

\bibitem{Sagan}
B. E. Sagan,
\emph{The Symmetric Group: Representations, Combinatorial Algorithms,
and Symmetric Functions}, second edition,
Graduate Texts in Mathematics 203, Springer, 2001.
\href{https://doi.org/10.1007/978-1-4757-6804-6}{doi:10.1007/978-1-4757-6804-6}.

\bibitem{StanleyFactorizations}
R. P. Stanley, \emph{Factorization of permutations into $n$-cycles},
Discrete Mathematics \textbf{37} (1981), 255--262.
\href{https://doi.org/10.1016/0012-365X(81)90224-7}{doi:10.1016/0012-365X(81)90224-7}.

\bibitem{StanleySurvey}
R. P. Stanley, \emph{Smith Normal Form in Combinatorics}, 2016.
\href{https://arxiv.org/abs/1602.00166}{arXiv:1602.00166}.

\bibitem{StanleyQuestion}
R. P. Stanley, \emph{Smith normal form of conjugacy class actions},
MathOverflow, July 1, 2019.
\url{https://mathoverflow.net/questions/335226}.
\end{thebibliography}
\end{document}
